\subsection{The actual harmonic heat comparison}
\label{actual-theta-heat-comparison}

Retain the original Host, CM type $\Psi$, even $n$, and $r$ with
$4r<n$. All Fourier measures, doubled coordinates and Gaussian
determinant factors are the prescribed ones. For an admitted family
$\varphi=(\varphi_\Lambda)$ put
\[
 D(Z,z,z')=\sum_{[\Lambda]}w_M(\Lambda)
       \bigl(\theta_{\varphi_\Lambda}(Z,z,z')
                         -E_{\varphi_\Lambda}(Z,z,z')\bigr).
\]
Theorem~\ref{actual-hermitian-siegel-weil-theorem} gives
$D(Z,0,0)=0$. This is not a uniqueness assertion for arbitrary
functions of the Jacobi variables. The following result proves
precisely the derivatives required by the spanning argument.

\begin{theorem}\label{actual-theta-heat-theorem}
Let $(P_j)$ be any finite family of the original harmonic polynomials
whose coefficient pairing is invariant under the column action of
$\prod_{\tau\in\Psi}U(n)$. For every pair of original frames
$x_1,x_2$ in the period domain,
\begin{equation}\label{actual-theta-heat-vanishing}
 \sum_j P_j(\partial_{z_1})\overline{P_j}(\partial_w)
 \left[
 K_0(x_1,x_2;z_1,w)
 D\bigl(Z_{12},z_{12}(z_1,w),z'_{12}(z_1,w)\bigr)
 \right]_{z_1=w=0}=0.
\end{equation}
Here the derivatives use the original letter directions;
$Z_{12},z_{12},z'_{12}$ and $K_0$ are literally the original doubled
period, linear Jacobi maps and Gaussian, including its determinant
factor. This supplies the original invariant harmonic heat comparison
for the actual class-sum difference, with no new analytic hypothesis
and no scalar normalization change.
The source target is the original contracted-jet theorem and its use in
\texttt{KernelIdentity}, not an assertion of full Jacobi equality
\cite{actual-sw-source}.
\end{theorem}

\begin{proof}
\textbf{The actual heat equation.}
The theta and Eisenstein expressions are holomorphic jointly in
$Z,z,z'$ on the original domain, and satisfy
\begin{equation}\label{actual-theta-heat-equation}
 \Delta_{\tau ab}D:=
 \sum_{k=1}^{n}\partial_{z_{\tau ka}}\partial_{z'_{\tau kb}}D
       =2\pi i\,\partial_{Z_{\tau ab}}D.
\end{equation}
For theta series, differentiate their defining exponential; its
quadratic coefficient is the product of the two linear coefficients.
On a compact subset of the period domain, its positive imaginary
part is bounded below. A Gaussian therefore dominates every fixed
derivative, uniformly on bounded Jacobi sets.

For an Eisenstein term write $A=CZ+D_0$ and $T=A^{-1}C$.
Its Jacobi factor is exactly
\[
 \det(A)^{-n}
 \exp\!\left(-2\pi i\sum_k z'_kTz_k\right).
\]
The identities $d\log\det A=\operatorname{tr}(T\,dZ)$ and
$dT=-T(dZ)T$ prove \eqref{actual-theta-heat-equation}, including
the term $-2\pi i nT_{ba}$ from the $n$ diagonal contractions.
This calculation retains the original exponent $n$.
On a compact set with $\operatorname{Im}Z\ge\delta I$ and bounded
coordinates, the unitary block relations give
\[
 \|T\|\le\delta^{-1},\qquad
 |\det(CZ+D_0)|\ge c_K|\det(CiI+D_0)|,
 \quad c_K>0.
\]
For the second inequality, normalize $(C,D_0)$ by
$(CC^*+D_0D_0^*)^{-1/2}$ and use compactness; positivity of
$\operatorname{Im}Z$ excludes a zero determinant. The same block
relations give the first inequality. Consequently the absolute
value of each term on that compact set is bounded by a fixed
multiple of its absolute value at $(iI,0,0)$. The absolutely
convergent series in the range $4r<n$ is thus locally normally
convergent, including the Jacobi variables. Cauchy's formula
justifies all derivatives. The finite weighted class sum preserves
the equation. No uniform estimate over different arithmetic settings
has been used.

\textbf{A direct invariant-polynomial argument.}
For $n\ge r$, every polynomial $R(X,Y)$ on
$\operatorname{Mat}_{n,r}(\mathbb C)\times
\operatorname{Mat}_{r,n}(\mathbb C)$ invariant under
$(X,Y)\mapsto(uX,Yu^{-1})$, $u\in U(n)$, is a polynomial in the
entries of $YX$. Here is a proof in exactly the needed range.
The unitary group is Zariski dense in $\mathrm{GL}_n(\mathbb C)$:
analytic continuation along exponentials of skew-Hermitian matrices,
whose complex span is all matrices, gives the same polynomial
identities on the general linear group. On the dense open set
$\det(YX)\ne0$, send $X$ to $X_0=(I_r,0)^t$. Write the resulting
$Y=(T,B)$. The stabilizer
$\left(\begin{smallmatrix}I_r&T^{-1}B\\0&I\end{smallmatrix}\right)$
then sends $Y$ to $(T,0)$. Thus
\[
 R(X,Y)=s(YX),\qquad s(T):=R(X_0,(T,0)).
\]
The right side is a polynomial, so the identity extends across the
complement of that dense open set. Applying this argument one
embedding at a time proves the product-group version. When $r=0$
the assertion is simply that a polynomial in no variables is constant.

\textbf{Absorb the Gaussian without losing its coefficients.}
Write $y=(y^-,y^+)$ for the two independent doubled Jacobi matrices.
Their column action is $(y^-,y^+)\mapsto
(u y^-,\overline u y^+)$, so the invariant pairings are
$q_{\tau ab}=\sum_k y^-_{\tau ka}y^+_{\tau kb}$.
The original maps $z_{12},z'_{12}$ act only on the row coordinates;
their two complementary row blocks therefore intertwine precisely
this action. The invariant pairing of the $P_j$ becomes an invariant
constant-coefficient differential symbol $Q$ after this change of
variables. The Gaussian is a nonzero frame-dependent scalar times
the exponential of a linear combination of the $q_{\tau ab}$.

Leibniz's rule gives a finite polynomial $R$ characterized by
\[
 Q(\partial_y)\bigl(K(y)F(y)\bigr)\big|_{y=0}
       =R(\partial_y)F(0).
\]
Explicitly its symbol is
$R(\xi)=Q(\partial_y)[K(y)e^{\langle y,\xi\rangle}]_{y=0}$.
This expression retains every Gaussian coefficient, factorial and
the constant determinant multiplier. Invariance of $Q$, $K$ and
evaluation at zero makes $R$ invariant for the dual column action.
The preceding argument expresses it as a polynomial $s$ in the
pairing symbols. Hence the left side of
\eqref{actual-theta-heat-vanishing} is
\[
 \left.s\bigl(\Delta_{\tau ab}\bigr)D(Z,z,z')
                                      \right|_{Z=Z_{12},\ z=z'=0}
 =\left.s\bigl(2\pi i\partial_{Z_{\tau ab}}\bigr)
                                      D(Z,0,0)\right|_{Z=Z_{12}}=0.
\]
All derivatives commute on the open holomorphic domain. The constant-term
contribution also vanishes because $D(Z,0,0)=0$. This proves the
literal contracted-jet identity, rather than falsely asserting that
all Jacobi derivatives vanish.
\end{proof}

The result applies to the original doubled family whenever that family
has been shown admitted. In the kernel identity the doubled-test
scalar $\kappa\ne0$, the class mass and the Gaussian reproducing
scalar $\rho>0$ remain unchanged. Neither the finite Schwartz system
nor a uniform global period norm follows from the heat equation alone.
